\section{Prédiction QH1a falsifiable}
\label{sec:prediction_qh1a}

Les Sections~\ref{sec:champ_connexion}--\ref{sec:decouplage_crt}
établissent la position structurelle de la PT par rapport au
débat fondationnel : l'angle d'holonomie est primitif, les
objets classiques sont dérivés, les corrélations entre canaux
modulaires sont arithmétiques et invariantes par la géométrie
émergente. Cette position n'est pas qu'interprétative : elle
implique une prédiction expérimentale précise, falsifiable sur
plusieurs plateformes existantes, et structuralement
distinguable de toute prédiction d'une physique standard sur
variété courbe.

\subsection{Énoncé de la prédiction}
\label{ssec:qh1a_enonce}

\begin{prediction}[QH1a : invariance géométrique de la MI]
\label{pred:qh1a}
Dans un système atomique multi-électronique sondé par
spectroscopie CRT (canaux $p = 3, 5, 7$), l'information mutuelle
$I(n^{(p)}\,;\,n^{(q)})$ entre nombres d'occupation modulaires :
\begin{enumerate}[label=(\roman*),nosep]
  \item est non-nulle et donnée par la formule fermée du
  Corollaire~\ref{cor:closed_form_fr} (test arithmétique) ;
  \item est invariante sous toute transformation géométrique
  du système (déplacement nucléaire, déformation de la maille
  cristalline, séparation des détecteurs) (test géométrique).
\end{enumerate}
\end{prediction}

La condition (i) est le contenu quantitatif de
Corollaire~\ref{cor:closed_form_fr} : la MI prend les valeurs
$\mathcal{I}_{3,5} = 0{,}138$ bits et
$\mathcal{I}_{3,7} = 0{,}283$ bits (valeurs corpus), prédictibles
au $\pm 20\%$ près par la formule fermée à un seul paramètre
entier $k_{p,q}$ par paire. Toute mesure qui ne reproduit pas ces
valeurs (à la précision statistique de l'expérience) invalide la
prédiction.

La condition (ii) est le contenu structurel de
Théorème~\ref{thm:invariance_geom} : la MI doit rester constante
sous toute reparamétrisation de la géométrie spatiale. Toute
mesure qui détecte une dépendance de la MI en la séparation des
détecteurs, la distance internucléaire, ou tout autre paramètre
géométrique, invalide la prédiction.

\subsection{Critère de falsification}
\label{ssec:qh1a_falsification}

Soit $G_0, G_1$ deux configurations géométriques distinctes d'un
même système quantique, et soient
$\mathcal{I}_{p,q}(G_0), \mathcal{I}_{p,q}(G_1)$ les valeurs de
l'information mutuelle mesurée sous chacune. Définissons
\begin{equation}
  \Delta\mathcal{I}_{p,q}
  \;:=\;
  \mathcal{I}_{p,q}(G_1) - \mathcal{I}_{p,q}(G_0)
  \label{eq:Delta_I_def}
\end{equation}
et soit $\epsilon_{\rm stat}$ l'incertitude statistique de
$\Delta\mathcal{I}_{p,q}$ telle qu'évaluée par la procédure de
mesure. La PT prédit
\begin{equation}
  \bigl|\Delta\mathcal{I}_{p,q}\bigr| \;<\; \epsilon_{\rm stat}
  \qquad \text{pour toute paire } (G_0, G_1).
  \label{eq:QH1a_predict}
\end{equation}
\emph{Critère de falsification.} Si une expérience mesure
$|\Delta\mathcal{I}_{p,q}| > 5 \epsilon_{\rm stat}$ pour une
paire de configurations géométriques distinctes, la prédiction
QH1a est falsifiée, et avec elle le théorème
d'invariance~\ref{thm:invariance_geom}. Une telle falsification
serait également une réfutation directe de la PT, puisque
$\theta_p$ et $\mathcal{A}_p$ sont des objets structuraux de la
théorie et ne peuvent être ajustés pour absorber un signal
géométrique.

\subsection{Plateformes expérimentales candidates}
\label{ssec:qh1a_plateformes}

Trois plateformes contemporaines offrent l'accès aux nombres
d'occupation modulaires $n^{(p)}$ avec contrôle indépendant de la
géométrie spatiale.

\paragraph{(a) Atomes froids dans réseaux optiques.}
Un réseau optique à espacement variable permet de scanner la
distance entre puits de potentiel adjacents tout en mesurant la
distribution d'occupation modulaire des niveaux internes par
spectroscopie de fluorescence. La géométrie est contrôlée par la
longueur d'onde du laser de piégeage ; les nombres d'occupation
$n^{(p)}$ sont accessibles par mesure résonante au niveau du
résidu modulaire. Précision attendue : $\epsilon_{\rm stat}
\sim 10^{-3}$ bits sur $\mathcal{I}_{3,5}$ pour
$N \sim 10^4$ atomes piégés, soit une sensibilité largement
suffisante pour falsifier QH1a à $5\sigma$ en présence d'un
couplage géométrique.

\paragraph{(b) Ions piégés à séparations contrôlées.}
Une chaîne d'ions Yb$^+$ ou Be$^+$ permet de varier la séparation
inter-ion avec une précision micrométrique, tout en mesurant les
corrélations inter-canaux par tomographie d'état. La structure
hyperfine des niveaux internes fournit les nombres d'occupation
modulaires pour $p = 3$ (canal de couleur SU(3) effectif) et
$p = 5$ (canal de saveur dans les configurations à plusieurs
électrons de valence). Précision attendue : $\epsilon_{\rm stat}
\sim 10^{-4}$ bits, soit la plateforme la plus contraignante.

\paragraph{(c) Molécules diatomiques à distance internucléaire ajustable.}
Une molécule diatomique sondée par spectroscopie laser à
résolution sub-MHz permet de varier la distance internucléaire
$R$ via un champ électrique externe (effet Stark) et de mesurer
simultanément les nombres d'occupation modulaires des niveaux
électroniques. La signature recherchée est l'absence de
dépendance de $\mathcal{I}_{3,5}$ en $R$ sur une plage de
plusieurs angstroms. Précision attendue : $\epsilon_{\rm stat}
\sim 5 \times 10^{-3}$ bits, plus modeste mais structuralement
complémentaire (le couplage géométrique attendu d'une théorie
standard est ici dominé par la variation du potentiel
vibrationnel).

\subsection{Contraste avec la physique standard}
\label{ssec:qh1a_contraste}

Une théorie quantique standard sur variété courbe prédit
génériquement $|d\mathcal{I}_{p,q}/dG| \neq 0$. La raison
structurelle est que les champs de matière sont, dans le
formalisme standard, des sections de fibrés au-dessus de la
variété courbe, et leurs corrélations héritent de la structure
métrique de la base par les couplages gravitationnels minimaux.
L'effet est typiquement faible (suppression par $1/M_{\rm Pl}^2$
ou par le couplage gravitationnel local), mais il est
\emph{non-nul} et calculable.

La PT, au contraire, prédit $d\mathcal{I}_{p,q}/dG = 0$
identiquement. La raison est que les canaux modulaires
$\mathcal{A}_p$ sont définis avant la métrique $G$ et lui sont
mesurablement antérieurs. Aucune fonction de la métrique
n'apparaît dans la formule
fermée~\eqref{eq:closed_form_fr} ; la dérivée par rapport à $G$
est exactement zéro.

Ce contraste rend QH1a un test discriminant fort entre les deux
cadres. Une mesure positive ($|d\mathcal{I}/dG| > 5\epsilon$)
falsifie la PT et confirme le cadre standard ; une mesure nulle
($|d\mathcal{I}/dG| < \epsilon$) confirme la PT au niveau
structurel et exclut tout couplage géométrique minimal au-dessus
de la précision atteinte. Aucune des plateformes actuelles n'a
encore réalisé la mesure complète --- en particulier, aucune
n'a scanné la dépendance géométrique de la MI sur plusieurs ordres
de grandeur de séparation. C'est la principale lacune
expérimentale à combler pour soumettre la PT à un test
falsificationniste réel.
